\Needspace{6\baselineskip}
\begin{center}\hypertarget{native-229-article166}{}\pdfbookmark[2]{166.}{article-166}166.\end{center}
Ad determinationem primam orbitae penitus adhuc incognitae e quatuor
observationibus necessario eiusmodi positiones adhibendae erunt, quae motum
heliocentricum non nimis magnum complectuntur: alioquin enim careremus
subsidiis ad approximationem primam commode formandam. Methodus tamen ea
quam statim trademus extensione tam lata gaudet, ut absque haesitatione
observationes motum heliocentricum $30$ vel $40$ graduum complectentes in
usum vocare liceat, si modo distantiae a Sole non nimis inaequales fuerint:
quoties eligendi copia datur, temporum intervalla inter primam et secundam,
secundam et tertiam, tertiam et quartam ab aequalitate parum recedentia accipere
iuvabit. Sed hoc quoque respectu anxietate nimia haud opus erit, uti exemplum
subnexum monstrabit, ubi temporum intervalla sunt $48$, $55$ et $59$ dierum,
motusque heliocentricus ultra $50^0$.

Porro solutio nostra requirit, ut completae sint observatio secunda et tertia,
adeoque latitudines vel declinationes in observationibus extremis \SourceLeaf{230}negligantur.
Supra quidem monuimus, praecisionis maioris gratia plerumque praestare,
si elementa duabus observationibus extremis completis, atque intermediarum
longitudinibus vel ascensionibus rectis accomodentur: attamen in prima orbitae
determinatione huic lucro renuntiavisse haud poenitebit, quum approximatio
expeditissima longe maioris momenti sit, iacturamque illam, quae praecipue
tantum in longitudinem nodi atque inclinationem orbitae cadit, elementaque
reliqua vix sensibiliter afficiat, postea facile explere liceat.

Brevitatis caussa methodi expositionem ita adornabimus, ut omnes locos ad
eclipticam referamus, adeoque quatuor longitudines cum duabus latitudinibus
datas esse supponemus: attamen quoniam in formulis nostris ad terrae latitudinem
quoque respicietur, sponte ad eum casum transferri poterunt, ubi aequator
tamquam planum fundamentale accipitur, si modo ascensiones rectae ac declinationes
in locum longitudinum et latitudinum substituuntur.

Ceterum respectu nutationis, praecessionis et parallaxis, nec non aberrationis,
omnia quae in Sectione praec. exposuimus etiam hic valent: nisi itaque
distantiae approximatae a terra aliunde iam innotuerunt, ut respectu aberrationis
methodum I art. 118 in usum vocare liceat, loca observata initio tantum ab
aberratione fixarum purgabuntur, temporaque corrigentur, quamprimum inter
calculi decursum distantiarum determinatio approximata in potestatem venit,
ut infra clarius elucebit.
